\chapter{Tensors} \label{ch:tensors}

\section{Tensors, Contraction, and Tensor Network}

We present the following straightforward definition of tensors, which is a rephrasal of Penrose's original definition \cite{penrose1971applications}.

\begin{definition}
Given $n\in \N$, let $[n]$ denote the set $\{0, 1, \cdots, n-1\}$. 
A tensor of rank $k \in \N$ is a function $T : [n_0] \times [n_1] \times \cdots [n_{k-1}] \rightarrow \C$.
\end{definition}

By covention, we shall write the value $T(i_0, i_1, \cdots, i_{k-1})$ as $T_{i_0 i_1 \cdots i_{k-1}}$, where each $i$ is called an index.
Of course, the output of $T$ can be $\R$ or any other ring. 
We ordain it to be $\C$ in this dissertation.
$n_j$ is called the dimension of the $j$-th index, which will always be clear from the context.

Following the convention of Einstein summation, a tensor $T$ is denoted by a capital letter with its indices as free variables shown in the subscript.
The class of all tensors is denoted by the caligraphy letter $\T_{ijk \cdots}$.

Tensor is, plainly and unonstensibly, an arrangement of complex numbers with some orderings by the indices.
For example, a rank-one tensor is a vector entries, rank-two tensor is a matrix, and a rank-three tensor is an arrangement of matrices.

The following equation denotes a vector and a matrix, while figure \ref{fig:rank3-stack} depicts the rank-three tensor as a stack of $n$ such matrix slices.
\begin{equation*}
	T_i =
	\begin{pmatrix}
		T_0 \\ T_1 \\ \vdots \\ T_{n-1}
	\end{pmatrix},
	\qquad
	T_{ij} =
	\begin{pmatrix}
		T_{00} & T_{01} & \cdots & T_{0,n-1} \\
		T_{10} & T_{11} & \cdots & T_{1,n-1} \\
		\vdots & \vdots & \ddots & \vdots \\
		T_{n-1,0} & T_{n-1,1} & \cdots & T_{n-1,n-1}
	\end{pmatrix}.
\end{equation*}

\begin{figure}[ht]
\centering
% A rank-three tensor $T_{ijk}$ regarded as a stack of $n$ matrix slices.
% Drawn back-to-front so the $i = 0$ slice is frontmost.
\begin{tikzpicture}[>=Stealth, line join=round, line cap=round]
    \def\s{2.4}     % matrix edge length
    \def\dx{0.55}   % per-slice depth offset (x)
    \def\dy{0.55}   % per-slice depth offset (y)

    % --- stack of n matrix slices, drawn back-to-front (i = 0 is frontmost) ---
    \foreach \k in {3,2,1,0} {
        \pgfmathsetmacro{\X}{\dx*\k}
        \pgfmathsetmacro{\Y}{\dy*\k}
        \pgfmathtruncatemacro{\idx}{\k}
        \begin{scope}[shift={({\X},{\Y})}]
            \fill[blue!8!white] (0,0) rectangle (\s,\s);
            \draw[thick] (0,0) rectangle (\s,\s);
            \node[font=\footnotesize, anchor=south east] at (0.04,\s) {$i=\idx$};
        \end{scope}
    }

    % internal grid on the front slice (i = 0)
    \foreach \i in {1,2,3} {
        \draw[black!30] (\i*0.6,0) -- (\i*0.6,\s);
        \draw[black!30] (0,\i*0.6) -- (\s,\i*0.6);
    }
    \node[font=\scriptsize] at (0.3,0.3) {$T_{000}$};
    \node[font=\scriptsize] at (0.9,0.3) {$T_{001}$};
    \node[font=\scriptsize] at (0.3,0.9) {$T_{010}$};

    % depth arrow indicating the n layers along the index i
    \draw[<->, thick] (-0.35 + 3.5,-0.35) -- ({3*\dx+3.5},{3*\dy+0.2});
    \node[font=\footnotesize, rotate=45] at ({3*\dx+3},{1.5*\dy-0.15}) {$n$ slices};
\end{tikzpicture}

\caption{A rank-three tensor $T_{ijk}$ regarded as a stack of $n$ matrices.}
\label{fig:rank3-stack}
\end{figure}



There are at least seven natural operations for tensors.
\begin{observation}
	\begin{itemize} 
		\item \textbf{Addition} $\T + \T \rightarrow \T$. Let $T_{ijk \cdots}$ and $S_{ijk \cdots}$ be two tensors of the same rank. Then their sum is defined by $(T + S)_{ijk \cdots} = T_{ijk \cdots} + S_{ijk \cdots}$.
		\item \textbf{Scalar multiplication} Let $T_{ijk \cdots}$ be a tensor and $c \in \C$. Define $(cT)_{ijk \cdots} =c (T_{ijk \cdots}) $
		\item \textbf{Permutation.} Let $T_{ijk\cdots}$ be a tensor of rank $m$. For any permutation $\sigma \in S_m$, the permuted tensor is defined by $(\sigma \cdot T)_{i_{\sigma(1)} i_{\sigma(2)} \cdots i_{\sigma(k)}} = T_{i_1 i_2 \cdots i_k}$.
		\item \textbf{Flatten} $\T \rightarrow \T$. Let $T_{ijk \cdots}$ be a tensor. 
		Any of the two indices of $T$ can be flattened into one to create a new tensor.
		Say we would like to flatten the indices $i$ and $j$, of dimension $n_1, n_2$. Let $\alpha$ be an index of dimension $n_1n_2$, and we define the new tensor $S_{ak \cdots}$ as $S_{(n_2 i + j)k \cdots} = T_{ijk\cdots}$.
		\item \textbf{Unflatten}, which is the operation exactly opposite of flatten.
		\item \textbf{Outer product} $\T \otimes \T \rightarrow \T$ Let $T_{ijk \cdots}$ and $S_{lmn \cdots}$ be any two tensors.
		Their tensor product is defined by $(T \otimes S)_{ijklmn\cdots} = T_{ijk\cdots} S_{lmn\cdots}$. Outer product is sometimes called \textbf{tensor product} or \textbf{Kronecker product}.
		\item \textbf{Contraction.} Let $T_{ijk\cdots}$ and $S_{ijk \cdots}$ be tensor of rank $m$ and $n$. Their contraction over the index $i$ is a tensor of rank $m+n - 2$ defined by $(T \cdot S)_{jk\cdots lm\cdots} = \sum_{i=0}^{n-1} T_{ijk\cdots} S_{ilm\cdots}$.
	\end{itemize}
\end{observation}

Addition and scalar multiplication of tensors are clear enough. 
Permutation of indices is generalisation of matrix transpose. 
For an example, consider ranked-two tensor $T$, and define $S_{ij} = T_{ji}$

\begin{equation*}
	T_{ij} =
	\begin{pmatrix}
		T_{00} & T_{01} & \cdots  \\
		T_{10} & T_{11} & \cdots \\ 
		\vdots & \vdots & \ddots \\
	\end{pmatrix}, \qquad
	S_{ij} =
	\begin{pmatrix}
		T_{00} & T_{10} & \cdots  \\
		T_{01} & T_{11} & \cdots \\ 
		\vdots & \vdots & \ddots \\
	\end{pmatrix}
\end{equation*}

Outer product and flattenning are demonstrated by figure \ref{fig:outer-flatten}.

\begin{figure}[ht]
\centering
% (a) Outer product of a vector and a matrix yields a rank-three tensor.
% (b) Flattening two indices (i, j) of that tensor yields a matrix.
% A concrete 2x2x2 instance is drawn so that every visible entry is labelled.
\begin{tikzpicture}[>=Stealth, line join=round, line cap=round]

    % ==========================================================
    % Panel (a): Outer product  v_i \otimes M_{jk} = T_{ijk}
    % ==========================================================
    \begin{scope}

        % --- vector v (length 2) ---
        \fill[blue!8!white] (0,0) rectangle (0.9,1.8);
        \draw[thick] (0,0) rectangle (0.9,1.8);
        \draw[black!30] (0,0.9) -- (0.9,0.9);
        \node[font=\scriptsize] at (0.45,1.35) {$v_0$};
        \node[font=\scriptsize] at (0.45,0.45) {$v_1$};
        \node[font=\scriptsize] at (0.45,-0.35) {$v_i$};

        % --- otimes ---
        \node[font=\large] at (1.5,0.9) {$\otimes$};

        % --- matrix M (2x2) ---
        \fill[blue!8!white] (2.1,0) rectangle (3.9,1.8);
        \draw[thick] (2.1,0) rectangle (3.9,1.8);
        \draw[black!30] (3.0,0) -- (3.0,1.8);
        \draw[black!30] (2.1,0.9) -- (3.9,0.9);
        \node[font=\scriptsize] at (2.55,1.35) {$M_{00}$};
        \node[font=\scriptsize] at (3.45,1.35) {$M_{01}$};
        \node[font=\scriptsize] at (2.55,0.45) {$M_{10}$};
        \node[font=\scriptsize] at (3.45,0.45) {$M_{11}$};
        \node[font=\scriptsize] at (3.0,-0.35) {$M_{jk}$};

        % --- equals ---
        \node[font=\large] at (4.35,0.9) {$=$};

        % --- rank-three tensor T (2x2x2): back slice (i=1) then front slice (i=0) ---
        \fill[blue!8!white] (5.2,0.4) rectangle (7.0,2.2);
        \draw[thick] (5.2,0.4) rectangle (7.0,2.2);
        \fill[blue!8!white] (4.8,0) rectangle (6.6,1.8);
        \draw[thick] (4.8,0) rectangle (6.6,1.8);
        \draw[black!30] (5.7,0) -- (5.7,1.8);
        \draw[black!30] (4.8,0.9) -- (6.6,0.9);
		\node[font=\scriptsize] at (5.25,1.35) {$v_0M_{00}$};
		\node[font=\scriptsize] at (6.15,1.35) {$v_0M_{01}$};
	\node[font=\scriptsize] at (5.25,0.45) {$v_0M_{10}$};
	\node[font=\scriptsize] at (6.15,0.45) {$v_0M_{11}$};
        \node[font=\scriptsize, anchor=south east] at (4.8,1.8) {$i\!=\!0$};
        \node[font=\scriptsize, anchor=south east] at (5.2,2.2) {$i\!=\!1$};
        \node[font=\scriptsize] at (5.7,-0.35) {$T_{ijk}$};

        \node[font=\small\itshape] at (3.5,-0.95) {(a) Outer product};
    \end{scope}

    % ==========================================================
    % Panel (b): Flatten  T_{ijk}  ->  S_{\alpha k},  \alpha = 2i + j
    % ==========================================================
    \begin{scope}[shift={(0,-4.9)}]

        % --- rank-three tensor T (input, 2x2x2) ---
        \fill[blue!8!white] (0.4,0.4) rectangle (2.2,2.2);
        \draw[thick] (0.4,0.4) rectangle (2.2,2.2);
        \fill[blue!8!white] (0,0) rectangle (1.8,1.8);
        \draw[thick] (0,0) rectangle (1.8,1.8);
        \draw[black!30] (0.9,0) -- (0.9,1.8);
        \draw[black!30] (0,0.9) -- (1.8,0.9);
        \node[font=\scriptsize] at (0.45,1.35) {$T_{000}$};
        \node[font=\scriptsize] at (1.35,1.35) {$T_{001}$};
        \node[font=\scriptsize] at (0.45,0.45) {$T_{010}$};
        \node[font=\scriptsize] at (1.35,0.45) {$T_{011}$};
        \node[font=\scriptsize, anchor=south east] at (0,1.8) {$i\!=\!0$};
        \node[font=\scriptsize, anchor=south east] at (0.4,2.2) {$i\!=\!1$};
        \node[font=\scriptsize] at (0.9,-0.35) {$T_{ijk}$};

        % --- flatten arrow ---
        \draw[->, very thick] (2.6,0.9) -- (4.5,0.9);
        \node[font=\footnotesize] at (3.55,1.25) {flatten};
        \node[font=\scriptsize] at (3.55,0.55) {$\alpha = 2i + j$};

        % --- flattened matrix S (4 rows x 2 cols) ---
        \fill[blue!8!white] (4.8,-0.9) rectangle (6.6,2.7);
        \draw[thick] (4.8,-0.9) rectangle (6.6,2.7);
        \draw[black!30] (5.7,-0.9) -- (5.7,2.7);
        \draw[black!30] (4.8,1.8) -- (6.6,1.8);
        \draw[black!30] (4.8,0.9) -- (6.6,0.9);
        \draw[black!30] (4.8,0.0) -- (6.6,0.0);
        \node[font=\scriptsize] at (5.25, 2.25) {$S_{00}$};
        \node[font=\scriptsize] at (6.15, 2.25) {$S_{01}$};
        \node[font=\scriptsize] at (5.25, 1.35) {$S_{10}$};
        \node[font=\scriptsize] at (6.15, 1.35) {$S_{11}$};
        \node[font=\scriptsize] at (5.25, 0.45) {$S_{20}$};
        \node[font=\scriptsize] at (6.15, 0.45) {$S_{21}$};
        \node[font=\scriptsize] at (5.25,-0.45) {$S_{30}$};
        \node[font=\scriptsize] at (6.15,-0.45) {$S_{31}$};
        \node[font=\scriptsize] at (5.7,-1.25) {$S_{\alpha k}$};

        \node[font=\small\itshape] at (3.3,-1.8) {(b) Flatten};
    \end{scope}

\end{tikzpicture}

\caption{(a) The outer product of a vector $v_i$ and a matrix $M_{jk}$ yields the rank-three tensor $T_{ijk}$. (b) Flattening the indices $i$ and $j$ of $T_{ijk}$ into a single index $\alpha = n_2\, i + j$ yields the matrix $S_{\alpha k}$.}
\label{fig:outer-flatten}
\end{figure}

Usually, outer product and flattenning are used together to convert tensor products of two arbitrary tensors into a single matrix, for example
\begin{equation}
	A \otimes \vec{v} = 
	\begin{pmatrix}
		a_{00} & a_{01}\\
		a_{10} & a_{11}
	\end{pmatrix}
	\begin{pmatrix}
		v_0 \\ v_1	
	\end{pmatrix}
	= 
	\begin{pmatrix}
		a_{00} v_0 & a_{01} v_0 \\
		a_{10} v_0 & a_{11} v_0 \\
		a_{00} v_1 & a_{01} v_1 \\
		a_{10} v_1 & a_{11} v_1
	\end{pmatrix}
\end{equation}

Contraction, however, is a more subtle operation and the key of this thesis, which is preserved for the next section.

\begin{remark}
It is worth noting that, in general, there is an intrinsic notion of input and output for in linear algebra. 
For example, in vector space $V$ a vector is an element of $V$. Its dual is a function $V \rightarrow \C$, and a matrix is a function $V \rightarrow V$.
In physics, the input and output of tensor are differentiated by the upper and lower indices. 
For example, $\T^k_l$ denotes the class of tensor $ \underbrace{V \otimes V \otimes \cdots \otimes V}_{k} \otimes \underbrace{V^* \otimes V^* \otimes \cdots \otimes V^*}_{l} \cong \underbrace{V \otimes  \cdot \otimes V}_{k} \rightarrow \underbrace{V \otimes \cdots \otimes V}_{l}$,
where $V$ is a vector space and $V^*$ is its dual.
While differentiating input and ouputs seems an intuitive idea, it actually is unweldy in practice, especially considering that for finite-dimensional vector space any $\T^k_l \cong \T^{k'}_{l'}$ provided $k + l = k' + l'$.
\end{remark}
